\documentclass[11pt,a4paper]{article}

\usepackage[utf8]{inputenc}
\usepackage[
  left=1.05in,
  right=1.05in,
  top=1in,
  bottom=1in,
  headheight=14pt
]{geometry}
\usepackage{amsmath,amssymb,amsthm,mathtools}
\usepackage{bm}
\usepackage{enumitem}
\usepackage{booktabs,array}
\usepackage{microtype}
\usepackage{xcolor}
\usepackage[breakable,skins]{tcolorbox}
\usepackage[hidelinks]{hyperref}
\usepackage{fancyhdr}

\linespread{1.25}
\setlength{\parindent}{0pt}
\setlength{\parskip}{1em plus 0.25em minus 0.1em}
\setlength{\abovedisplayskip}{1.3em plus 0.3em minus 0.2em}
\setlength{\belowdisplayskip}{1.3em plus 0.3em minus 0.2em}
\setlength{\abovedisplayshortskip}{1em plus 0.2em}
\setlength{\belowdisplayshortskip}{1.1em plus 0.2em}
\setlength{\jot}{10pt}
\setlength{\emergencystretch}{2em}
\raggedbottom

\numberwithin{equation}{section}

\newtheoremstyle{spacedplain}
  {18pt}{18pt}{\itshape}{}{\bfseries}{.}{0.5em}{}
\newtheoremstyle{spaceddefinition}
  {18pt}{18pt}{\normalfont}{}{\bfseries}{.}{0.5em}{}
\newtheoremstyle{spacedremark}
  {18pt}{18pt}{\normalfont}{}{\itshape}{.}{0.5em}{}

\theoremstyle{spacedplain}
\newtheorem{theorem}{Theorem}[section]
\newtheorem{proposition}[theorem]{Proposition}

\theoremstyle{spaceddefinition}
\newtheorem{definition}[theorem]{Definition}

\theoremstyle{spacedremark}
\newtheorem{remark}[theorem]{Remark}

\setlist[itemize]{leftmargin=*,itemsep=7pt,topsep=8pt,parsep=3pt}
\setlist[enumerate]{leftmargin=*,itemsep=7pt,topsep=8pt,parsep=3pt}

% Plain black frame, used only for the most important results.
\newtcolorbox{keybox}{
  enhanced,
  colback=white,colframe=black,
  boxrule=0.6pt,sharp corners,
  left=10pt,right=10pt,top=4pt,bottom=6pt,
  before skip=16pt,after skip=16pt
}

\pagestyle{fancy}
\fancyhf{}
\fancyhead[L]{MTL712}
\fancyhead[R]{Lecture 35}
\fancyfoot[C]{\thepage}
\renewcommand{\headrulewidth}{0.4pt}
\renewcommand{\footrulewidth}{0pt}

% Run-in headings inside proofs and sections, set on their own line.
\newcommand{\step}[1]{\par\addvspace{0.6em}\noindent\textbf{#1}\par\nobreak}

\newcommand{\pd}[2]{\frac{\partial #1}{\partial #2}}
\newcommand{\pdd}[2]{\frac{\partial^2 #1}{\partial #2^2}}

\title{\textbf{MTL712: Computational Methods for Differential\\Equations}\\[6pt]
\large Lecture 35: ADI Methods and Three-Level Schemes}
\author{\normalsize Typeset lecture notes}
\date{\normalsize 30 October 2026}

\begin{document}

\maketitle
\thispagestyle{plain}
{\linespread{1.0}\selectfont\setlength{\parskip}{0pt}\tableofcontents}

%=====================================================================
\section{Alternating Direction Implicit (ADI) Methods}
%=====================================================================

Consider the two-dimensional heat equation
\[
  \pd ut=Lu,
  \qquad
  L=D_1^2+D_2^2,
  \qquad
  D_1=\pd{}{x_1},\quad D_2=\pd{}{x_2},
\]
that is,
\[
  \pd ut=\pdd u{x_1}+\pdd u{x_2},
\]
with proper initial and boundary conditions.  As in one dimension, the
exact solution satisfies
\[
  u_{l,m}^{n+1}=\exp(kL)\,u_{l,m}^n .
\]

\subsection{The implicit (Crank--Nicolson type) scheme}

Split the exponential symmetrically about \(t_{n+1/2}\):
\[
  \exp\bigl(-\tfrac12kL\bigr)u_{l,m}^{n+1}
  =\exp\bigl(\tfrac12kL\bigr)u_{l,m}^n,
\]
that is,
\begin{equation}\label{eq:adi1}
  \exp\bigl(-\tfrac12k(D_1^2+D_2^2)\bigr)u_{l,m}^{n+1}
  =\exp\bigl(\tfrac12k(D_1^2+D_2^2)\bigr)u_{l,m}^n .
\end{equation}
The derivative operators are expanded in central differences:
\begin{align*}
  D_1^2&=\frac1{h^2}\Bigl[\delta_{x_1}^2-\frac1{12}\delta_{x_1}^4
         +\frac1{90}\delta_{x_1}^6-\cdots\Bigr],\\
  D_2^2&=\frac1{h^2}\Bigl[\delta_{x_2}^2-\frac1{12}\delta_{x_2}^4
         +\frac1{90}\delta_{x_2}^6-\cdots\Bigr].
\end{align*}
Retaining only the second-order central differences, and using
\(r=k/h^2\), equation \eqref{eq:adi1} becomes
\[
  \exp\bigl(-\tfrac12r\delta_{x_1}^2\bigr)
  \exp\bigl(-\tfrac12r\delta_{x_2}^2\bigr)U_{l,m}^{n+1}
  =\exp\bigl(\tfrac12r\delta_{x_1}^2\bigr)
   \exp\bigl(\tfrac12r\delta_{x_2}^2\bigr)U_{l,m}^n .
\]
Keeping the first two terms of each exponential gives the following
scheme.

\begin{keybox}
\begin{proposition}\label{prop:adi}
The two-dimensional analogue of the Crank--Nicolson scheme is
\begin{equation}\label{eq:adi2}
  \bigl(1-\tfrac12r\delta_{x_1}^2\bigr)
  \bigl(1-\tfrac12r\delta_{x_2}^2\bigr)U_{l,m}^{n+1}
  =\bigl(1+\tfrac12r\delta_{x_1}^2\bigr)
   \bigl(1+\tfrac12r\delta_{x_2}^2\bigr)U_{l,m}^n .
\end{equation}
\end{proposition}
\end{keybox}

Recall that the one-dimensional Crank--Nicolson scheme is
\[
  \bigl(1-\tfrac12r\delta_x^2\bigr)U_m^{n+1}
  =\bigl(1+\tfrac12r\delta_x^2\bigr)U_m^n ;
\]
\eqref{eq:adi2} is its analogue in two dimensions.

\subsection{Splitting}

Split \eqref{eq:adi2} into two formulas, each implicit in only one
direction:
\begin{equation}\label{eq:split}
\begin{aligned}
  \bigl(1-\tfrac12r\delta_{x_1}^2\bigr)U_{l,m}^{n+1*}
  &=\bigl(1+\tfrac12r\delta_{x_2}^2\bigr)U_{l,m}^n,\\
  \bigl(1-\tfrac12r\delta_{x_2}^2\bigr)U_{l,m}^{n+1}
  &=\bigl(1+\tfrac12r\delta_{x_1}^2\bigr)U_{l,m}^{n+1*}.
\end{aligned}
\end{equation}
Here \(U^{n+1*}\) is an intermediate value.

Eliminating \(U_{l,m}^{n+1*}\) gives back the combined formula
\eqref{eq:adi2}, under the assumption that \(\delta_{x_1}^2\) and
\(\delta_{x_2}^2\) commute.  Indeed, multiplying the second equation of
\eqref{eq:split} by \(\bigl(1-\tfrac12r\delta_{x_1}^2\bigr)\) and
using the first,
\begin{align*}
  \bigl(1-\tfrac12r\delta_{x_1}^2\bigr)
  \bigl(1-\tfrac12r\delta_{x_2}^2\bigr)U_{l,m}^{n+1}
  &=\bigl(1+\tfrac12r\delta_{x_1}^2\bigr)
    \bigl(1-\tfrac12r\delta_{x_1}^2\bigr)U_{l,m}^{n+1*}\\
  &=\bigl(1+\tfrac12r\delta_{x_1}^2\bigr)
    \bigl(1+\tfrac12r\delta_{x_2}^2\bigr)U_{l,m}^n .
\end{align*}

\subsection{The Peaceman--Rachford ADI method}

Taking the intermediate value to be the value at the half time level,
\[
  U_{l,m}^{n+1*}=U_{l,m}^{n+\frac12},
\]
in \eqref{eq:split} gives the Peaceman--Rachford formula.

\begin{keybox}
\begin{definition}[Peaceman--Rachford ADI method]
\begin{align*}
  \bigl(1-\tfrac12r\delta_{x_1}^2\bigr)U_{l,m}^{n+\frac12}
  &=\bigl(1+\tfrac12r\delta_{x_2}^2\bigr)U_{l,m}^n,\\
  \bigl(1-\tfrac12r\delta_{x_2}^2\bigr)U_{l,m}^{n+1}
  &=\bigl(1+\tfrac12r\delta_{x_1}^2\bigr)U_{l,m}^{n+\frac12}.
\end{align*}
\end{definition}
\end{keybox}

\begin{itemize}
  \item In the \emph{first step} we advance from \(t_n\) to
    \(t_{n+1/2}\), using an implicit formula for
    \(\partial^2u/\partial x_1^2\) and an explicit formula for
    \(\partial^2u/\partial x_2^2\).
  \item In the \emph{second step} we advance from \(t_{n+1/2}\) to
    \(t_{n+1}\), using an explicit difference for
    \(\partial^2u/\partial x_1^2\) and an implicit difference for
    \(\partial^2u/\partial x_2^2\).
\end{itemize}

%=====================================================================
\section{A Three-Level Scheme}
%=====================================================================

A three-level scheme involves three time levels, \(t_{n-1}\), \(t_n\)
and \(t_{n+1}\).  Consider
\[
  \pd ut=\pdd ux,
\]
with
\[
  \text{IC: } u(x,0)=f(x),
  \qquad
  \text{BC: } u(0,t)=u(1,t)=0,\quad t>0
\]
(or other appropriate initial and boundary conditions).

\step{Difference scheme}
Use a central difference in time as well as in space:
\[
  \frac{U_m^{n+1}-U_m^{n-1}}{2k}
  =\frac1{h^2}\bigl(U_{m+1}^n-2U_m^n+U_{m-1}^n\bigr).
\]
With \(r=k/h^2\) this is
\[
  U_m^{n+1}=U_m^{n-1}+2r\bigl(U_{m+1}^n-2U_m^n+U_{m-1}^n\bigr),
\]
that is,
\begin{equation}\label{eq:richardson}
  U_m^{n+1}=2r\,\delta_x^2U_m^n+U_m^{n-1}.
\end{equation}
(This three-level scheme is known as Richardson's scheme.)

\step{Reduction to a two-level system}
Define \(V_m^n=U_m^{n-1}\).  Then \eqref{eq:richardson} becomes
\[
  U_m^{n+1}=2r\,\delta_x^2U_m^n+V_m^n,
  \qquad
  V_m^{n+1}=U_m^n .
\]

\subsection{Stability analysis (von Neumann)}

Write the system in matrix--vector form:
\[
  \begin{bmatrix}U\\ V\end{bmatrix}_m^{n+1}
  =\begin{bmatrix}2r\delta_x^2&1\\ 1&0\end{bmatrix}
   \begin{bmatrix}U\\ V\end{bmatrix}_m^{n}.
\]
With \(W=\begin{bmatrix}U\\ V\end{bmatrix}\) we have
\[
  W_m^{n+1}=\begin{bmatrix}2r\delta_x^2&1\\ 1&0\end{bmatrix}W_m^n .
\]
Let
\[
  W_m^n=W_0^n\,e^{i\beta mh},
\]
where \(W_0^n\) is a constant vector; this is a typical Fourier term.
In components,
\[
  U_m^n=U_0^n\,e^{i\beta mh},
  \qquad
  V_m^n=V_0^n\,e^{i\beta mh}.
\]
Substituting,
\begin{align*}
  U_0^{n+1}e^{i\beta mh}
  &=2r\Bigl[U_0^ne^{i\beta(m+1)h}-2U_0^ne^{i\beta mh}
    +U_0^ne^{i\beta(m-1)h}\Bigr]+V_0^ne^{i\beta mh},\\
  V_0^{n+1}e^{i\beta mh}&=U_0^ne^{i\beta mh}.
\end{align*}
Since \(e^{i\beta(m+1)h}+e^{i\beta(m-1)h}=2\cos\beta h\,e^{i\beta mh}\),
\begin{align*}
  U_0^{n+1}e^{i\beta mh}
  &=2r\bigl(2\cos\beta h-2\bigr)U_0^ne^{i\beta mh}+V_0^ne^{i\beta mh},\\
  V_0^{n+1}e^{i\beta mh}&=U_0^ne^{i\beta mh}.
\end{align*}
Therefore
\[
  W_m^{n+1}=\begin{bmatrix}2r(2\cos\beta h-2)&1\\ 1&0\end{bmatrix}W_m^n,
\]
and, using \(2\cos\beta h-2=-4\sin^2\frac{\beta h}2\),
\[
  W_m^{n+1}=\underbrace{\begin{bmatrix}-8r\sin^2\frac{\beta h}2&1\\[2pt] 1&0\end{bmatrix}}_{\text{amplification matrix}}W_m^n .
\]

\begin{definition}[Von Neumann condition for a two-level system]
The von Neumann necessary condition for stability of a two-level
system is
\[
  \max_{i}|\lambda_i|\le1,
\]
where the \(\lambda_i\) are the eigenvalues of the amplification
matrix.
\end{definition}

\begin{keybox}
\begin{theorem}\label{thm:richardson}
The three-level scheme \eqref{eq:richardson} is unconditionally
unstable.
\end{theorem}
\end{keybox}

\begin{proof}
Let \(\lambda_i\), \(i=1,2\), be the eigenvalues of the amplification
matrix.  They satisfy
\[
  \lambda_i^2+8r\lambda_i\sin^2\frac{\beta h}2-1=0,
\]
so
\begin{align*}
  \lambda_{1,2}
  &=\frac{-8r\sin^2\frac{\beta h}2
     \pm\bigl(64r^2\sin^4\frac{\beta h}2+4\bigr)^{1/2}}{2}\\
  &=-4r\sin^2\frac{\beta h}2
     \pm\Bigl(1+16r^2\sin^4\frac{\beta h}2\Bigr)^{1/2}.
\end{align*}
That is,
\begin{align*}
  \lambda_1&=-4r\sin^2\frac{\beta h}2
    +\Bigl(1+16r^2\sin^4\frac{\beta h}2\Bigr)^{1/2},\\
  \lambda_2&=-4r\sin^2\frac{\beta h}2
    -\Bigl(1+16r^2\sin^4\frac{\beta h}2\Bigr)^{1/2}.
\end{align*}
For every \(r>0\) and every \(\beta\) with
\(\sin\frac{\beta h}2\ne0\), both terms of \(\lambda_2\) are negative
and the square root exceeds \(1\), so
\[
  |\lambda_2|
  =4r\sin^2\frac{\beta h}2+\Bigl(1+16r^2\sin^4\frac{\beta h}2\Bigr)^{1/2}
  >1 .
\]
The von Neumann condition is therefore violated by \(\lambda_2\) for
all values of \(r\), so the method is unconditionally unstable.
\end{proof}

%=====================================================================
\section{The Du Fort--Frankel Scheme}
%=====================================================================

Consider again
\[
  \pd ut=\pdd ux
\]
with appropriate initial and boundary conditions, and \(r=k/h^2\).

Start from the explicit three-level scheme \eqref{eq:richardson}:
\[
  U_m^{n+1}=2r\bigl(U_{m+1}^n+U_{m-1}^n-2U_m^n\bigr)+U_m^{n-1}.
\]
Replace \(U_m^n\) by the average of its neighbours in time,
\[
  U_m^n\ \longrightarrow\ \frac{U_m^{n+1}+U_m^{n-1}}2 .
\]
We get
\[
  U_m^{n+1}=2r\bigl(U_{m+1}^n+U_{m-1}^n\bigr)
  -2r\bigl(U_m^{n+1}+U_m^{n-1}\bigr)+U_m^{n-1},
\]
so
\[
  (1+2r)\,U_m^{n+1}
  =2r\bigl(U_{m+1}^n+U_{m-1}^n\bigr)+(1-2r)\,U_m^{n-1}.
\]

\begin{keybox}
\begin{definition}[Du Fort--Frankel scheme]
\[
  U_m^{n+1}
  =\frac{2r}{1+2r}\bigl(U_{m+1}^n+U_{m-1}^n\bigr)
   +\frac{1-2r}{1+2r}\,U_m^{n-1}.
\]
This scheme is unconditionally stable.
\end{definition}
\end{keybox}

\end{document}
